% ============================================================
%  sec-1-overview.tex  —  3.1 Nemotron 3 Nano Overview
% ============================================================

\section{Nemotron 3 Nano Overview}
\label{sec:nemotron-overview}

Nemotron 3 Nano is an open, efficient large language model developed by
NVIDIA and released in 2025. It belongs to the Nemotron 3 model family —
a suite of models spanning multiple capability tiers — and is positioned
as a high-efficiency alternative to models with significantly larger active
parameter counts. According to the repository root \texttt{README.md},
Nemotron 3 Nano \textit{``achieves better or on-par accuracy than competitive
models while having up to 3.3$\times$ higher inference throughput''} compared
to models such as GPT-OSS-20B and Qwen3-30B-A3B.

The model is designed for \textbf{agentic reasoning} workloads: multi-step,
multi-turn tasks requiring tool invocation, long-context comprehension, and
structured output generation. Its primary use cases, as documented in the
repository, include agentic pipelines with tool calling, extended context
document analysis (up to 1M tokens), multilingual tasks across 15+ languages,
and reasoning tasks with on/off reasoning-mode control.

% ── 3.1.1 ────────────────────────────────────────────────────
\subsection{Architecture}
\label{subsec:nemotron-architecture}

Nemotron 3 Nano implements a \term{Hybrid Mamba-Transformer} architecture
combined with a \term{sparse Mixture-of-Experts} (\acrshort{moe}) design.
This constitutes a departure from dense transformer-only architectures:
the hybrid design interleaves Mamba state-space model layers with standard
transformer attention layers, enabling efficient long-sequence processing
without the quadratic complexity of full self-attention across all layers.

The sparse \acrshort{moe} component activates only a subset of expert
networks per token during each forward pass, decoupling total parameter
count from per-token compute cost. \tabref{nemotron-arch-specs} documents
the architectural parameters extracted from the repository
\texttt{README.md}.

\begin{table}[htbp]
  \centering
  \caption{Nemotron 3 Nano architectural specifications}
  \label{tab:nemotron-arch-specs}
  \begin{tabularx}{0.88\textwidth}{lX}
    \toprule
    \tableheader{Parameter & Value} \\
    \midrule
    Total parameters         & 31.6B \\
    \rowcolor{TableRowOdd}
    Active parameters        & 3.6B (per forward pass) \\
    Architecture type        & Hybrid Mamba-Transformer with sparse MoE \\
    \rowcolor{TableRowOdd}
    Number of layers         & 32 \\
    Model dimension          & 3,008 \\
    \rowcolor{TableRowOdd}
    Total routable experts   & 128 \\
    Activated experts        & 6 per token per MoE layer \\
    \rowcolor{TableRowOdd}
    Shared experts           & 2 \\
    Pretraining tokens       & 25 trillion \\
    \rowcolor{TableRowOdd}
    Context length           & Up to 1M tokens \\
    Training stages          & 3 (Pretrain $\to$ SFT $\to$ RL) \\
    \bottomrule
  \end{tabularx}
\end{table}

The 31.6B/3.6B parameter ratio means that each forward pass activates
approximately 11.4\% of the total parameters. This is the efficiency
mechanism that enables inference throughput 3.3$\times$ higher than dense
models of similar capability.

% ── 3.1.2 ────────────────────────────────────────────────────
\subsection{Model Capabilities}
\label{subsec:nemotron-capabilities}

The following capabilities are documented in the repository
\texttt{README.md} and constitute the target requirements for the full
training pipeline described in subsequent sections.

% ── 3.1.3 ────────────────────────────────────────────────────
\subsection{Model Family Positioning}
\label{subsec:model-family}

The Nemotron 3 family, as referenced in the repository, comprises three
tiers: \textbf{Nano} (31.6B total / 3.6B active, optimised for throughput
and edge/cloud inference), \textbf{Super} (mid-tier capability/cost
balance), and \textbf{Ultra} (maximum capability, targeting frontier
benchmark performance). Nano's design rationale is explicit in the
repository: it trades total parameter count for inference efficiency by
exploiting sparse expert activation, making it deployable on hardware
configurations that cannot sustain dense models of equivalent benchmark
performance. The training pipeline — documented across four stage
directories — is shared in architecture across the family, with
stage-specific configurations controlling model dimension and parallelism
settings per tier.
